% ============================================================
%  Divisibilidade — Apostila Premium | PratiqueJá
%  Engine: XeLaTeX
% ============================================================
\documentclass[12pt]{article}
\usepackage{../../pratiqueja}

\setsubject{Divisibilidade}

\begin{document}
\pagenumbering{arabic}
\setstretch{1.2}
\color{bodytext}

\heroheader{Divisibilidade}%
           {Regras, Divisores e Aplicações}%
           {Matemática · Intermediário}

\setlength{\columnsep}{18pt}
\setlength{\columnseprule}{0.5pt}
\def\columnseprulecolor{\color{exborder}}

\begin{multicols}{2}

%─────────────────────────────────────────────────────────────
\TW{Def}{badgepurple}{Definição}{Divisibilidade}

\regra{Dizemos que $a$ é \textbf{divisível} por $b$ quando existe
$q \in \mathbb{Z}$ tal que $a = b \cdot q$, ou seja, o \textbf{resto
da divisão é zero}. Chamamos $b$ de \textbf{divisor} de $a$ e $a$
de \textbf{múltiplo} de $b$.}

%─────────────────────────────────────────────────────────────
\TW{2}{darkblue}{Divisibilidade por 2}{Números pares}

\regra{O \textbf{último dígito} deve ser \textbf{0, 2, 4, 6 ou 8}.}

\exemp{%
  13\hl{8} → último dígito 8 → divisível \ok\\
  47\hl{5} → último dígito 5 → não divisível \nok%
}

%─────────────────────────────────────────────────────────────
\TW{3}{darkblue}{Divisibilidade por 3}{Soma dos algarismos}

\regra{A \textbf{soma dos algarismos} deve ser múltipla de \textbf{3}.}

\exemp{%
  132 → 1+3+2=\hl{6} → múltiplo de 3 \ok\\
  531 → 5+3+1=\hl{9} → múltiplo de 3 \ok\\
  124 → 1+2+4=\hl{7} → não múltiplo de 3 \nok%
}

%─────────────────────────────────────────────────────────────
\TW{4}{darkblue}{Divisibilidade por 4}{Dois últimos dígitos}

\regra{Os \textbf{2 últimos dígitos} devem formar número divisível
por 4. Os demais algarismos não influenciam.}

\exemp{%
  5\hl{24} → 24 é múltiplo de 4 \ok\\
  13\hl{36} → 36 é múltiplo de 4 \ok\\
  5\hl{14} → 14 não é múltiplo de 4 \nok%
}

%─────────────────────────────────────────────────────────────
\TW{5}{darkblue}{Divisibilidade por 5}{Último dígito 0 ou 5}

\regra{O \textbf{último dígito} deve ser \textbf{0} ou \textbf{5}.}

\exemp{%
  32\hl{5} → termina em 5 \ok\quad
  14\hl{0} → termina em 0 \ok\\
  78\hl{3} → termina em 3 → não divisível \nok%
}

%─────────────────────────────────────────────────────────────
\TW{6}{darkblue}{Divisibilidade por 6}{Divisível por 2 e por 3}

\regra{Deve ser \textbf{par} e a soma dos algarismos deve ser
\textbf{múltipla de 3} — ambas as condições simultaneamente.}

\exemp{%
  654: par e 6+5+4=\hl{15} (÷3) \ok\\
  312: par e 3+1+2=\hl{6} (÷3) \ok\\
  125: ímpar → não divisível por 6 \nok%
}
\obs{Basta uma condição falhar para o número não ser divisível por 6.}

%─────────────────────────────────────────────────────────────
\TW{7}{darkblue}{Divisibilidade por 7}{Regra do duplo}

\regra{Multiplique o \textbf{último dígito por 2} e subtraia do
número sem esse dígito. Repita até concluir.}

\exemp{%
  16\hl{1} → 1·2=2 → 16−2=\hl{14} \ok\\
  109\hl{9} → 9·2=18 → 109−18=91 → 9−2=\hl{7} \ok%
}
\obs{Para números pequenos, a divisão direta costuma ser mais prática.}

%─────────────────────────────────────────────────────────────
\TW{8}{darkblue}{Divisibilidade por 8}{Três últimos dígitos}

\regra{Os \textbf{3 últimos dígitos} devem formar número divisível
por 8. Os demais não influenciam.}

\exemp{%
  7\hl{128} → 128 é múltiplo de 8 \ok\\
  5\hl{256} → 256 é múltiplo de 8 \ok\\
  3\hl{100} → 100 não é múltiplo de 8 \nok%
}

%─────────────────────────────────────────────────────────────
\TW{9}{darkblue}{Divisibilidade por 9}{Soma dos algarismos}

\regra{A \textbf{soma dos algarismos} deve ser múltipla de \textbf{9}.}

\exemp{%
  432 → 4+3+2=\hl{9} \ok\quad
  124 → 1+2+4=\hl{7} \nok\\
  6561 → 6+5+6+1=\hl{18} → múltiplo de 9 \ok%
}
\obs{Todo número divisível por 9 também é divisível por 3.}

%─────────────────────────────────────────────────────────────
\TW{10}{darkblue}{Divisibilidade por 10}{Último dígito 0}

\regra{O \textbf{último dígito deve ser 0}. Como $10=2{\times}5$,
equivale a ser divisível por 2 e por 5 simultaneamente.}

\exemp{%
  15\hl{0} → termina em 0 \ok\quad
  487\hl{0} → termina em 0 \ok\\
  32\hl{5} → termina em 5 → divisível por 5, mas não por 10 \nok%
}

%─────────────────────────────────────────────────────────────
\TW{11}{badgegreen}{Divisibilidade por 11}{Soma alternada dos algarismos}

\regra{A \textbf{soma alternada} dos algarismos (alternando $+$ e
$-$ da esquerda para a direita) deve ser múltipla de 11 ou zero.}

\exemp{%
  253 → +2−5+3=\hl{0} → múltiplo de 11 \ok\\
  1452 → +1−4+5−2=\hl{0} \ok\\
  123 → +1−2+3=\hl{2} → não divisível \nok%
}

%─────────────────────────────────────────────────────────────
\TW{Div}{badgegreen}{Divisores de um Número}{Como encontrar todos os divisores naturais}

Para cada fator primo da decomposição, multiplique \textbf{todos os
divisores já encontrados} por esse fator e acrescente os novos.

\SH{Exemplo — divisores de 60}

$60 = 2^2 \cdot 3^1 \cdot 5^1$

\vspace{3pt}
\noindent
$\begin{array}[t]{r|l|l}
  & & \textcolor{darkblue}{1} \\
60 & \textcolor{darkblue}{2} & \textcolor{darkblue}{2} \\
30 & \textcolor{darkblue}{2} & \textcolor{darkblue}{4} \\
15 & \textcolor{darkblue}{3} & \textcolor{darkblue}{3,\ 6,\ 12} \\
 5 & \textcolor{darkblue}{5} & \textcolor{darkblue}{5,\ 10,\ 20,\ 15,\ 30,\ 60} \\
\hline
 1 & & \\
\end{array}$

\SH{Todos os divisores de 60}

\noindent
\chip{1}\chip{2}\chip{3}\chip{4}%
\chip{5}\chip{6}\chip{10}\chip{12}%
\chip{15}\chip{20}\chip{30}\chip{60}

\SH{Quantidade de divisores}

Se $n = p_1^{a_1} \cdot p_2^{a_2} \cdots$, o total de divisores é:

\formula{$d(n) = (a_1+1)\cdot(a_2+1)\cdots$}

\exemp{%
  $60 = 2^2{\cdot}3^1{\cdot}5^1
  \Rightarrow d(60)=(3)(2)(2)=\hl{12}$ divisores\\
  $36 = 2^2{\cdot}3^2
  \Rightarrow d(36)=(3)(3)=\hl{9}$ divisores%
}

\end{multicols}

\end{document}
